\section{帰無仮説検定のシミュレーション}
\subsection{授業内容}
\subsubsection{科目の中でのこのコマの位置づけ}
標本統計量が確率変数であるなら，「二群の差」といった標本の群間差も確率的に変動するし，これに意味があるかないかといった判断も確率的になる。
この判断が誤ったものになる水準をコントロールしようというのが帰無仮説検定の考え方である。このことに注意しておけば，帰無仮説検定の手続きも「米印を探す単調作業」でないことは理解できるだろう。
改めて帰無仮説検定の手順を確認するとともに，t検定とシミュレーションを通じて検定の精度と意義の理解を深める。なお\textcite{kosugi2023}の第5章を参考にする。
\subsubsection{コマ主題細目}
\begin{description}
	\item[帰無仮説検定の論理] 帰無仮説と対立仮説，判断の手続などを再確認する。また，用語としてのタイプ1エラー，タイプ2エラー，検定力，および帰無仮説が従う分布(帰無分布)についても解説を加える。
		\begin{flushright}
			$\to$\textcite{kosugi2023}のP.187--191.
		\end{flushright}
	\item[Rによる検定の実際] 二群の平均値差の検定を例にとる。データを乱数から生成し，どのような手順で検定が行われているか，そのアルゴリズムを復習する。
		\begin{flushright}
			$\to$\textcite{kosugi2023}のP.191--196.
		\end{flushright}
	\item[タイプ1エラー確率のシミュレーション] 心理学の研究実践において，タイプ1エラー，タイプ2エラーがどれぐらいだったか，検定力がどれぐらいあったかを知る術はない。しかしシミュレーションすることによって，何がどの程度起こりうるかを体験することができる。これをもって，自らの研究実践に対する統計的指標の意義を考える一助にしてもらいたい。また，t検定においてはWelchの方法がデフォルトで用いられるが，等分散性の仮定が成立しない場合はどのような問題が起きるかも確認できる。
		\begin{flushright}
			$\to$\textcite{kosugi2023}のP.196--203.
		\end{flushright}
\end{description}
\subsubsection{キーワード}
\begin{itemize}
	\item 帰無仮説と対立仮説
	\item 帰無分布
	\item サンプルサイズ，タイプ1エラー，タイプ2エラー，検定力
	\item 等分散性の仮定
\end{itemize}
\subsection{授業情報}
\paragraph{コマの展開方法}
講義/遠隔/演習
\subsubsection{予習・復習課題}
\paragraph{予習} 一年時の「データ解析基礎」のテキストを復習しておくこと。第18--20講が本時に対応する。
\paragraph{復習} サンプルサイズや二群それぞれの分散の値をさまざまに変化させて，シミュレーションの結果がどのように変わるかいろいろ試してみよう。